\documentclass[border=12pt]{standalone}
\usepackage[utf8]{inputenc}
\usepackage[T1]{fontenc}
\usepackage[american, siunitx]{circuitikz}

\sisetup{output-decimal-marker={,}}

\ctikzset{
    resistors/scale=0.8
}

\begin{document}
\begin{circuitikz}[font=\sffamily\small]

    % =========================================================================
    % 1. BLOCO ESQUERDO: ARDUINO UNO R3 - MESTRE (U1)
    % =========================================================================
    \draw[thick, fill=gray!5] (-7.2, -3.2) rectangle (-2.4, 3.2);
    \node[align=center, font=\sffamily\bfseries] at (-5.4, 0.4) {U1\\Arduino\\Uno R3};
    \node[font=\sffamily\scriptsize, text=gray!70!black] at (-5.4, -0.7) {(Mestre)};

    % Alimentacao de U1
    \draw (-4.8, 3.2) -- (-4.8, 4.0) node[vcc]{+5V};
    \node[below, font=\ttfamily\footnotesize] at (-4.8, 3.2) {5V};

    \draw (-4.8, -3.2) -- (-4.8, -4.0) node[ground]{};
    \node[above, font=\ttfamily\footnotesize] at (-4.8, -3.2) {GND};

    % Pinos de comunicacao e controle de direcao de U1
    \node[left, font=\ttfamily\scriptsize] at (-2.4, 1.6) {D3 (SW RX)};
    \node[left, font=\ttfamily\scriptsize] at (-2.4, 0.0) {D2 (DIR)};
    \node[left, font=\ttfamily\scriptsize] at (-2.4, -1.6) {D4 (SW TX)};

    % =========================================================================
    % 2. TRANSCEPTOR RS-485 DO MESTRE: MAX485 (U2)
    % =========================================================================
    \draw[thick, fill=gray!5] (0.8, -3.2) rectangle (5.4, 3.2);
    \node[align=center, font=\sffamily\bfseries] at (3.2, 0.3) {U2\\MAX485};
    \node[font=\sffamily\scriptsize, text=gray!70!black] at (3.2, -0.5) {(RS-485)};

    % Alimentacao de U2
    \draw (3.1, 3.2) -- (3.1, 4.0) node[vcc]{+5V};
    \node[below, font=\ttfamily\scriptsize] at (3.1, 3.2) {VCC (8)};

    \draw (3.1, -3.2) -- (3.1, -4.0) node[ground]{};
    \node[above, font=\ttfamily\scriptsize] at (3.1, -3.2) {GND (5)};

    % Conexoes U1 -> U2
    \draw (-2.4, 1.6) -- (0.8, 1.6);
    \node[right, font=\ttfamily\scriptsize] at (0.8, 1.6) {RO (1)};

    \draw (-2.4, 0.0) -- (-0.2, 0.0) node[circ]{};
    \draw (0.8, 0.4) -- (-0.2, 0.4) -- (-0.2, -0.4) -- (0.8, -0.4);
    \node[right, font=\ttfamily\scriptsize] at (0.8, 0.4) {$\overline{\text{RE}}$ (2)};
    \node[right, font=\ttfamily\scriptsize] at (0.8, -0.4) {DE (3)};

    \draw (-2.4, -1.6) -- (0.8, -1.6);
    \node[right, font=\ttfamily\scriptsize] at (0.8, -1.6) {DI (4)};

    % Pinos do barramento em U2
    \node[left, font=\ttfamily\scriptsize] at (5.4, 1.2) {A (6)};
    \node[left, font=\ttfamily\scriptsize] at (5.4, -1.2) {B (7)};

    % =========================================================================
    % 3. BARRAMENTO RS-485 (A, B) E RESISTORES DE TERMINACAO (R1, R2 = 120R)
    % =========================================================================
    \draw (5.4, 1.2) -- (11.4, 1.2) node[midway, above, font=\ttfamily\footnotesize] {A};
    \draw (5.4, -1.2) -- (11.4, -1.2) node[midway, below, font=\ttfamily\footnotesize] {B};

    % Terminacao lado Mestre (R1) e lado Escravo (R2)
    \draw (7.0, 1.2) to[R, l^=R1, a_=120<\ohm>, *-*] (7.0, -1.2);
    \draw (9.8, 1.2) to[R, l^=R2, a_=120<\ohm>, *-*] (9.8, -1.2);

    % =========================================================================
    % 4. TRANSCEPTOR RS-485 DO ESCRAVO: MAX485 (U3) - ESPELHADO
    % =========================================================================
    \draw[thick, fill=gray!5] (11.4, -3.2) rectangle (16.0, 3.2);
    \node[align=center, font=\sffamily\bfseries] at (13.6, 0.3) {U3\\MAX485};
    \node[font=\sffamily\scriptsize, text=gray!70!black] at (13.6, -0.5) {(RS-485)};

    % Alimentacao de U3
    \draw (13.7, 3.2) -- (13.7, 4.0) node[vcc]{+5V};
    \node[below, font=\ttfamily\scriptsize] at (13.7, 3.2) {VCC (8)};

    \draw (13.7, -3.2) -- (13.7, -4.0) node[ground]{};
    \node[above, font=\ttfamily\scriptsize] at (13.7, -3.2) {GND (5)};

    % Pinos do barramento em U3
    \node[right, font=\ttfamily\scriptsize] at (11.4, 1.2) {A (6)};
    \node[right, font=\ttfamily\scriptsize] at (11.4, -1.2) {B (7)};

    % Conexoes U3 -> U4
    \draw (16.0, 1.6) -- (19.2, 1.6);
    \node[left, font=\ttfamily\scriptsize] at (16.0, 1.6) {RO (1)};

    \draw (19.2, 0.0) -- (17.0, 0.0) node[circ]{};
    \draw (16.0, 0.4) -- (17.0, 0.4) -- (17.0, -0.4) -- (16.0, -0.4);
    \node[left, font=\ttfamily\scriptsize] at (16.0, 0.4) {$\overline{\text{RE}}$ (2)};
    \node[left, font=\ttfamily\scriptsize] at (16.0, -0.4) {DE (3)};

    \draw (16.0, -1.6) -- (19.2, -1.6);
    \node[left, font=\ttfamily\scriptsize] at (16.0, -1.6) {DI (4)};

    % =========================================================================
    % 5. BLOCO DIREITO: ARDUINO UNO R3 - ESCRAVO (U4)
    % =========================================================================
    \draw[thick, fill=gray!5] (19.2, -3.2) rectangle (24.0, 3.2);
    \node[align=center, font=\sffamily\bfseries] at (22.2, 0.4) {U4\\Arduino\\Uno R3};
    \node[font=\sffamily\scriptsize, text=gray!70!black] at (22.2, -0.7) {(Escravo)};

    % Alimentacao de U4
    \draw (21.6, 3.2) -- (21.6, 4.0) node[vcc]{+5V};
    \node[below, font=\ttfamily\footnotesize] at (21.6, 3.2) {5V};

    \draw (21.6, -3.2) -- (21.6, -4.0) node[ground]{};
    \node[above, font=\ttfamily\footnotesize] at (21.6, -3.2) {GND};

    % Pinos de comunicacao e controle de direcao de U4
    \node[right, font=\ttfamily\scriptsize] at (19.2, 1.6) {D3 (SW RX)};
    \node[right, font=\ttfamily\scriptsize] at (19.2, 0.0) {D2 (DIR)};
    \node[right, font=\ttfamily\scriptsize] at (19.2, -1.6) {D4 (SW TX)};

\end{circuitikz}
\end{document}